\section{The single-line functional on the marked square}\label{app:weights}
The second boundary mean compares joint insertion with its first-order approximation. Inserting \(e\)
changes the weight of \(f\) because the insertion deletes the row and column of the inserted square.
We order \(e,f\) by increasing row and retain the two crossing indicators \(X,Y\). Here \(M_e\) is the
single rook minor obtained by deleting the row and column of \(e\). Primes denote evaluation on that
minor. Let \(t=p'_f/p_f-1\), \(G_{ef}=\mathsf G(e,f)\) and \(D=\chi_e+\chi_f-G_{ef}\). The exact
product of the insertion ratios gives
\[
 e^{u}=\frac{\pi^T_{ef}}{p_ep_f}
      =(1+\chi_e)(1+\chi'_f)(1+t)\exp(r_e+r'_f).
\]
The rook increment estimate is \(r_i=-\mathsf E_i+O(j^{-3})\), also on the single minor. With
\(t=-G_{ef}+O(j^{-2})\) and \(\chi'_f-\chi_f=O(j^{-2})\), the physical remainder divided by \(p_ep_f\)
is \(\eta+O(j^{-3})\), where
\begin{align*}
 \eta={}&(t+G_{ef})+(\chi'_f-\chi_f)+\chi_e\chi_f+t(\chi_e+\chi_f)\\
        &-\mathsf E_e-\mathsf E'_f+(p_e+p_f)G_{ef}\\
     ={}&B_{\rm phys}+D^2/2+O(j^{-3}),\\
 B_{\rm phys}={}&u-D+(p_e+p_f)G_{ef}.
\end{align*}
These seven terms apply in every state and orientation. The square term \(D^2/2\) requires only the
linear term of the insertion. To identify the remainder \(B_{\rm phys}\), we must compare the action
and prefactor on the mask and its double minor. To do this within one family we remove only \(e,f\)
and retain their four lines. Both margins of square \(i\) become \(1-x_i\), and all other margins
equal one. Let \(F_-(x),\ell_-(x)\) be the stationary action and the rook correction of this punctured
mask. Set \(R(x)=F_-(x)+2\sum_i(1-x_i)\log(1-x_i)\) and \(A(x)=\ell_-(x)+\sum_i\log(1-x_i)\). Both
regularized expressions extend to \(x=(1,1)\), where the visible system is the double minor and the
total mass of its Hessian preserves the correct size constant. Let \(p=(p_e,p_f)^\top\),
\(\mathbf1=(1,1)^\top\) and \(D_p=\operatorname{diag}(p_i(1-p_i))\). For the punctured margin Hessian
\(T_-\) and the incidence \(a_i\), we have \(R''_{il}=2\delta_{il}/(1-x_i)-a_i^\top T_-^+a_l\).
Reinserting the squares adds \(\sum_i[p_i(1-p_i)/(1-2p_i)]a_i a_i^\top\). The determinant lemma on the
complement of the gauge and stationarity give
\begin{align*}
 u={}&R(\mathbf1)-R(p)-(\mathbf1-p)^\top\nabla R(p)\\
     &+A(\mathbf1)-A(p)+\tfrac12\log\det(I-R''(p)D_p).
\end{align*}
Let \(\partial=\partial_e+\partial_f\), \(H=R''(0)\) and \(\Delta C=C(\mathbf1)-C(0)\).
Differentiating the stationary equation with respect to the two weights gives the response
\((D_p^{-1}-R''(p))^{-1}\). We expand the determinant and the responses of \(\chi_e,\chi_f,G_{ef}\).
The terms linear in \(p_i H_{il}\) cancel by \(\mathbf1^\top Hp=\sum_i p_iH_{ii}+(p_e+p_f)H_{ef}\).
This leaves
\[
 B_{\rm phys}=\Delta R-\partial R(0)-\partial^2 R(0)/2+\Delta A-\partial A(0)+O_b(j^{-3}).
\]
To obtain the coefficients, we add two copies of opposite signs to each marked line. On the positive
copy \(\lambda_i=1-x_i\), and on the ordinary part and the negative copies \(\lambda_i=1\). With
auxiliary mass \(s\), the counting measure is \(M=a+s\rho\) and the margin measure is \(N=\lambda M\),
with \(\rho=0,-1,+1\) on the ordinary part, the negative copies and the positive copies, respectively.
At \(s=1/j\) the negative copies cancel the originals exactly. For the binary incidence \(Y_{il}\) and
its Sinkhorn factors~\cite{sk} \(W_{il}\), the normalized equations are
\[
 w_{il}(1-s\lambda_i\lambda_l w_{il})=W_{il}\exp(z_i+y_l),\qquad
 \int Y_{il}w_{il}\,dN_C(l)=\int Y_{il}w_{il}\,dN_R(i)=1 .
\]
Let \(R_q,A_q\) denote the coefficients of \(s^q\) in the regularized action and prefactor along
\(x_e=x_f=t\). In the action, the coefficient of \(s^0\) is affine and \(R_1\) has degree at most two
in \(t\). In the prefactor \(A_0\) is constant and \(A_1\) is affine, since at each marked endpoint
the two quadratic terms from the local factor and the margin diagonal have opposite signs and equal
magnitude. The rook gap and the kernel section bounds control the inverse of the margin Jacobian in
the gauge. The implicit equation on a fixed complex disk in \(s\), for \(|x_i|\le2\), and the Cauchy
estimates give uniform \(O_b(j^{-3})\) remainders for the regularized values and derivatives used
above. With primes denoting derivatives in \(t\), we obtain
\begin{align*}
 L_s&=A'_1(0)+\tfrac12R''_1(0),\\
 B_{2,s}&=\tfrac12\int_0^1(1-t)^2R'''_2(t)\,dt+\int_0^1(1-t)A''_2(t)\,dt .
\end{align*}
In \(L_s,B_{2,s}\) the subscript denotes the state of the crossings. The polynomials
\(R_1,R_2,A_1,A_2\) have degrees at most \(4,6,6,8\), so nine distinct nodes determine their
derivatives algebraically. Floating-point evaluation requires its own error control. To complete the
proof we must transfer \(L_s\) and \(B_{2,s}\) to the profile in each state, the first squared and the
second linearly.
